Le quattro tavole d'insieme del corpo mantengono i gruppi originali; qui
restano i dettagli dei singoli package e le collaborazioni ulteriori.
Le tavole estendono le viste selezionate nella sezione~\ref{sec:classi}.
Ogni titolo riporta l'identificativo stabile e il package descritto;
i tipi esterni servono soltanto a leggere una collaborazione e sono
riconoscibili dalle qualifiche. Il grigio identifica il contesto esterno
nelle tavole che lo dichiarano; in CL-08a e CL-09 distingue invece la
persistenza, come indicato nel titolo. Il colore non sostituisce mai
nome, stereotipo o notazione dell'arco.

I membri sono selezionati: \code{+}, \code{-}, \code{\#} e \texttt{\textasciitilde}
indicano rispettivamente visibilità pubblica, privata, protetta e di package;
\code{abstract}, \code{interface}, \code{enumeration} e \code{record}
qualificano la natura del tipo. Le firme abbreviate mantengono i nomi
dei metodi, omettendo parametri o tipi non necessari alla lettura.
I valori tra parentesi quadre indicano opzionalità o molteplicità;
\code{ID} e \code{chiave} segnalano riferimenti scalari.
Le note e i vincoli esplicativi non sono nuovi membri Java.

Una freccia continua aperta indica un riferimento navigabile;
una freccia tratteggiata aperta indica uso, parametro o costruzione.
Il triangolo vuoto distingue generalizzazione e realizzazione mediante
linea continua o tratteggiata. Le molteplicità sono annotate al bersaglio
e non dichiarano una proprietà esclusiva del collaboratore. Nessun rombo
implica un ciclo di vita non dimostrato. Gli attributi che nominano tipi
esterni esplicitano ulteriori collaborazioni senza ripetere tutti gli archi.

CL-01 copre avvio e composizione; CL-02--CL-05 i quattro package del dominio;
CL-06 servizi e record di lettura; CL-07 eventi e observer;
CL-08a--CL-08b contratti DAO e gestore transazionale; CL-09 le dieci
implementazioni PostgreSQL; CL-10 presentazione; CL-11 credenziali.
Le undici viste di package coprono tutti i 65 tipi principali correnti,
senza richiedere package intermedi vuoti. La selettività riguarda soprattutto
i membri, gli archi ripetuti e i tipi annidati: i record di lettura del
catalogo sono richiamati in nota e le callback transazionali sono esplicite.
La matrice package--tavola--tipi--sorgenti e il controllo delle relazioni
sono conservati nelle verifiche documentali del ticket 03.

\newcommand{\tavolaclassi}[4]{%
  \clearpage
  \begin{figure}[p]
    \centering
    \includegraphics[width=\linewidth,height=.88\textheight,keepaspectratio]{diagrammi/#1.drawio.pdf}
    \caption{#2 --- #3.}\label{#4}
  \end{figure}\clearpage}

\tavolaclassi{classi-avvio}{CL-01}{Avvio e composizione, \code{it.unifi.drivehub}}{fig:cl01}
\tavolaclassi{classi-utenti}{CL-02}{Identità, ruoli e conto, \code{domain.user}}{fig:cl02}
\tavolaclassi{classi-catalogo}{CL-03}{Catalogo e scelta della base, \code{domain.vehicle}}{fig:cl03}
\tavolaclassi{classi-sconti}{CL-04}{Sconti per ID, \code{domain.discount}}{fig:cl04}
\tavolaclassi{classi-operazioni}{CL-05}{Richieste, prenotazioni e registro, \code{domain.operation}}{fig:cl05}
\tavolaclassi{classi-servizi}{CL-06}{Servizi e collaborazioni, \code{service}}{fig:cl06}
\tavolaclassi{classi-eventi}{CL-07}{Eventi e observer sincroni, \code{event}}{fig:cl07}
\tavolaclassi{classi-dao}{CL-08a}{Contratti per risorsa, \code{dao}}{fig:cl08a}
\tavolaclassi{classi-transazioni}{CL-08b}{Gestore e callback transazionali, \code{dao}}{fig:cl08b}
\tavolaclassi{classi-postgres}{CL-09}{Realizzazioni PostgreSQL, \code{dao.postgres}}{fig:cl09}
\tavolaclassi{classi-presentazione}{CL-10}{Viste, dialoghi e caricamenti, \code{presentation.view}}{fig:cl10}
\tavolaclassi{classi-sicurezza}{CL-11}{Contratto e realizzazione per le password, \code{security}}{fig:cl11}
